% ===========================================================================
%  sviluppatore/03-architettura/01-esagonale.tex
% ===========================================================================
\section{L'architettura esagonale, in concreto}\label{sez:esagonale}

Il guscio a strati occupa \num{4546} righe su \num{119084}: meno del
\num{4}\,\% del Python del progetto. Non ha riscritto niente. Avvolge il
codice che c'era --- \pkgStrategie, \pkgReport, \pkgEsecuzione, \pkgLive,
\pkgDati{} --- e gli mette davanti un confine.

\figuraTikz{esagonale}
  {I cinque strati. L'unica regola è la direzione delle frecce.}
  {esagonale}

\subsection{La regola, e perché è una sola}

Si dipende solo verso l'interno. \pkgApplicazione{} non conosce nessun
concreto: conosce solo le porte. I concreti si scelgono in un file solo,
\file{composition/container.py}.

\begin{perche}{Perché il confine vale la pena}
L'ecosistema che \Lyra{} vuole essere --- applicazione scaricabile, sito
community, marketplace --- è fatto di \emph{tre client dello stesso motore}.
Finché il motore è una libreria che ciascun client importa a modo suo, ogni
client ne vede una parte diversa e ogni cambiamento ne rompe due. Rendendolo
un servizio con un contratto stabile, i client tornano a essere consumatori, e
l'inversione delle dipendenze non è più una buona pratica da ricordarsi: è
imposta dal confine di processo.
\end{perche}

\subsection{\pkgDominio: i tipi, e basta}

Quattrocentosettantotto righe di \cod{dataclass} congelate. Nessuna
dipendenza --- nemmeno pandas. \simbolo{StrategyInfo}, \simbolo{StreamInfo},
\simbolo{RunReport}, \simbolo{OrderRequest}, \simbolo{SecurityReport},
\simbolo{LiveRunnerConfig} e una trentina d'altri.

Sono il vocabolario che attraversa gli strati. Un caso d'uso restituisce un
\simbolo{StrategyInfo}, non un \simbolo{DataFrame} e non un
\simbolo{StrategyDescriptor} --- che è un oggetto del motore, con dentro una
\cod{factory} richiamabile.

\subsection{\pkgPorte: trentadue contratti}

Il file da leggere per primo in tutto il progetto, e sono
\num{239} righe. Ogni porta è un \cod{Protocol} marcato
\cod{@runtime\_checkable}.

\begin{codiceInterno}{python}{Due porte intere, non un estratto}{ports/\_\_init\_\_.py}
@runtime_checkable
class StrategyRepository(Protocol):
    """Sorgente delle strategie disponibili (mappate a StrategyInfo)."""
    def list(self) -> list[StrategyInfo]: ...
    def get(self, name: str) -> Optional[StrategyInfo]: ...


@runtime_checkable
class WalkForwardRunner(Protocol):
    """Esegue un walk-forward per una strategia. Ritorna l'exit code."""
    def run(self, strategy_name: str) -> int: ...
\end{codiceInterno}

\noindent
Sono \cod{Protocol}, non classi base astratte, e la differenza conta:

\begin{description}[leftmargin = 1.6em, style = nextline, itemsep = 4pt]
  \item[Tipizzazione strutturale] Un concreto soddisfa la porta
    \emph{senza ereditarla}. Non deve importare \pkgPorte, non deve
    dichiarare niente: le basta avere i metodi giusti. L'accoppiamento è il
    minimo possibile, e si può avvolgere codice preesistente senza
    toccarlo.
  \item[Porte piccole e separate] \simbolo{DataSource} $\neq$
    \simbolo{Broker} $\neq$ \simbolo{StrategyRepository} $\neq$
    \simbolo{Clock}. Un adattatore implementa quello che gli serve. È
    l'\emph{Interface Segregation} preso sul serio: una porta
    «\simbolo{Everything}» avrebbe costretto ogni finto a implementare
    trenta metodi per provarne uno.
  \item[Sostituibilità verificata] Nei test i casi d'uso girano con
    repository e runner finti. Non è una proprietà sperata: è il modo in cui
    le prove sono scritte.
\end{description}

\begin{nota}{La porta \simbolo{Broker} è ri-esportata, non riscritta}
\simbolo{live.broker.LiveBroker} era già una classe base astratta, cioè già
una porta, scritta prima del guscio. \pkgPorte{} la ri-esporta invece di
duplicarla:

\begin{codiceInterno}{python}{}{ports/\_\_init\_\_.py}
try:
    from live.broker import LiveBroker as Broker  # noqa: F401
except Exception:  # pragma: no cover
    Broker = object  # type: ignore
\end{codiceInterno}

\noindent
È la mossa corretta: due definizioni della stessa interfaccia divergono, e
quando divergono nessuno se ne accorge finché un broker non si comporta in un
modo in un posto e in un altro altrove.
\end{nota}

\subsection{\pkgApplicazione: i casi d'uso}

Ogni caso d'uso è una classe con una responsabilità, riceve le porte dal
costruttore e non conosce nessun concreto. Sono brevi --- spesso tre righe --- e
devono restarlo: il caso d'uso orchestra, non calcola.

\begin{codiceInterno}{python}{Tre casi d'uso, per intero}{application/\_\_init\_\_.py}
class ListStrategies:
    def __init__(self, repo: StrategyRepository) -> None:
        self._repo = repo

    def execute(self) -> list[StrategyInfo]:
        return self._repo.list()


class RunWalkForward:
    """Esegue un walk-forward e (se c'e' un sink) ne registra il report."""
    def __init__(self, runner: WalkForwardRunner,
                 report_sink: Optional[ReportSink] = None) -> None:
        self._runner = runner
        self._sink = report_sink
\end{codiceInterno}

\begin{attenzione}{Il calcolo non va qui}
La tentazione, aggiungendo una funzionalità, è di metterne «un pezzetto» nel
caso d'uso: un filtro, una media, una conversione. Il giorno in cui succede,
quel pezzetto esiste solo nel percorso che passa per l'API e non in quello
che passa per la dashboard, e i due danno numeri diversi. Il calcolo sta nel
motore; il caso d'uso lo chiama e basta.
\end{attenzione}

\subsection{\pkgComposizione: centonove righe, e sono quelle decisive}

\file{composition/container.py} è un \cod{dataclass} con venti campi --- un
caso d'uso ciascuno --- e una funzione che lo costruisce.

\begin{codiceInterno}{python}{L'unico punto di scelta dei concreti}{composition/container.py}
    # --- unico punto di scelta dei CONCRETI ---
    repo = RegistryStrategyRepository(Path(strategies_dir) if strategies_dir else None)
    weighting = EnsembleWeightingModeProvider()
    runner = InProcessWalkForwardRunner()
    registry = FileModelRegistryAdapter(model_dir)
    sink = FileReportSink(report_dir)
    ...
    return Container(
        list_strategies=ListStrategies(repo),
        get_strategy=GetStrategy(repo),
        run_walkforward=RunWalkForward(runner, sink),
        ...
    )
\end{codiceInterno}

\noindent
Per cambiare sorgente dati, broker o registro dei modelli si tocca questo
file, e nessun altro. Per far girare i casi d'uso nei test con dei finti si
costruisce un \simbolo{Container} diverso, senza toccare né i casi d'uso né
le porte.

\begin{pratica}{Come si aggiunge una capacità, nell'ordine giusto}
\begin{enumerate}
  \item Un tipo in \pkgDominio, se serve un dato nuovo che attraversa il
        confine.
  \item Una porta in \pkgPorte: che cosa deve saper fare il mondo.
  \item Un caso d'uso in \pkgApplicazione: tre righe, nessun calcolo.
  \item Un adattatore in \pkgAdattatori{} che avvolge il codice del motore
        che fa davvero la cosa.
  \item Una riga in \file{composition/container.py} che li collega.
  \item Un router in \pkgApi{} che lo espone, se serve via HTTP.
\end{enumerate}
Saltare il passo~\num{2} è l'errore tipico: si collega il caso d'uso
direttamente all'adattatore, «tanto è uno solo», e da quel momento il caso
d'uso non è più verificabile senza l'infrastruttura.
\end{pratica}
